% ============================================================
%  Guide d'Entretien — Analyse des Vulnérabilités Réseaux
%  IA & Automatisation de la Détection d'Intrusions
%  Projet Académique — Module : Intelligence Artificielle
% ============================================================
\documentclass[12pt,a4paper]{article}

\usepackage[utf8]{inputenc}
\usepackage[T1]{fontenc}
\usepackage[english]{babel}
\newcommand{\og}{\guillemotleft\,}
\newcommand{\fg}{\,\guillemotright}
\usepackage{geometry}
\usepackage{booktabs}
\usepackage{array}
\usepackage{enumitem}
\usepackage{fancyhdr}
\usepackage{titlesec}
\usepackage{hyperref}
\usepackage{parskip}
\usepackage{amssymb}

\geometry{margin=2.5cm, top=3cm, bottom=3cm}

% --- Column types ---
\newcolumntype{B}[1]{>{\bfseries\raggedright\arraybackslash}p{#1}}
\newcolumntype{L}[1]{>{\raggedright\arraybackslash}p{#1}}

% --- Section formatting ---
\titleformat{\section}
  {\normalfont\large\bfseries}{\thesection.}{0.6em}{}
\titlespacing{\section}{0pt}{14pt}{6pt}

% --- Header/Footer ---
\pagestyle{fancy}
\fancyhf{}
\fancyhead[L]{\small Guide d'Entretien --- Vuln\'{e}rabilit\'{e}s R\'{e}seaux}
\fancyhead[R]{\small Module : Intelligence Artificielle}
\fancyfoot[C]{\small Projet Acad\'{e}mique \quad Page \thepage}
\renewcommand{\headrulewidth}{0.4pt}

\hypersetup{colorlinks=false, hidelinks}

\setlength{\parskip}{6pt}
\setlength{\parindent}{0pt}
\setlength{\tabcolsep}{6pt}

\begin{document}

% ============================================================
%  HEADER BLOCK
% ============================================================
\begin{center}
  {\Large\bfseries Guide d'Entretien}\\[0.4em]
  {\large Analyse des Vuln\'{e}rabilit\'{e}s R\'{e}seaux}\\[0.2em]
  {\normalsize IA \& Automatisation de la D\'{e}tection d'Intrusions}\\[1em]
\end{center}

\begin{tabular}{@{}B{3cm} L{10cm}@{}}
  \hline
  Module    & Intelligence Artificielle \\
  \hline
  Statut    & Projet Acad\'{e}mique \\
  \hline
  \'{E}tudiant 1 & El Kartouti Anas \\
  \hline
  \'{E}tudiant 2 & Lekhbioui Adil \\
  \hline
  \'{E}tudiant 3 & Moulay Anas \\
  \hline
  \'{E}tudiant 4 & Hammoubel El Yazid \\
  \hline
  \'{E}tudiant 5 & Nafia Khadija \\
  \hline
\end{tabular}

\vspace{1.2em}

% ============================================================
%  1. PHASE D'INTRODUCTION
% ============================================================
\section{Phase d'Introduction (Mise en situation)}

\textbf{Q1.1 :} \og{} Pourriez-vous vous pr\'{e}senter bri\`{e}vement et d\'{e}crire votre r\^{o}le ainsi que le type
d'architecture r\'{e}seau (LAN, Cloud, r\'{e}seaux hybrides) que vous g\'{e}rez ou utilisez au
quotidien~? \fg{}

\bigskip

\textbf{Q1.2 :} \og{} De mani\`{e}re g\'{e}n\'{e}rale, quels sont les d\'{e}fis ou les incidents majeurs auxquels vous
\^{e}tes le plus fr\'{e}quemment confront\'{e} pour garantir la disponibilit\'{e} et la s\'{e}curit\'{e} de
cette infrastructure~? \fg{}

% ============================================================
%  2. PHASE DE DÉVELOPPEMENT
% ============================================================
\section{Phase de D\'{e}veloppement (Exploration du probl\`{e}me cible)}

\textbf{Q2.1 :} \og{} Pourriez-vous me d\'{e}crire en d\'{e}tail une exp\'{e}rience r\'{e}cente o\`{u} vous avez constat\'{e}
un comportement suspect, une anomalie de trafic ou une baisse de performance inexpliquée
sur le r\'{e}seau~? \fg{}

\bigskip

\textbf{Q2.2 :} \og{} Selon votre exp\'{e}rience, lorsque le r\'{e}seau subit ce type de perturbations de mani\`{e}re
furtive, quelles en sont les causes les plus courantes~? \fg{}

\bigskip

\textbf{Q2.3 :} \og{} Face \`{a} ces anomalies, comment faites-vous la diff\'{e}rence entre une simple panne de
composant, une inondation de trafic malveillante (DDoS/surcharge), ou une infection par
un malware interne~? \fg{}

\bigskip

\textbf{Q2.4 :} \og{} Les attaquants proc\`{e}dent souvent par des phases de reconnaissance avant de frapper.
Comment percevez-vous ou d\'{e}tectez-vous ces tentatives de cartographie silencieuse
(comme l'exploration de vos services ou les scans de ports) sur votre p\'{e}rim\`{e}tre~? \fg{}

% ============================================================
%  3. PHASE D'APPROFONDISSEMENT
% ============================================================
\section{Phase d'Approfondissement (Limites actuelles et transition vers l'innovation)}

\textbf{Q3.1 :} \og{} Actuellement, quelles proc\'{e}dures, logs ou outils mobilisez-vous pour isoler la cause
d'une intrusion ou d'un scan, et comment proc\'{e}dez-vous pour bloquer la menace~? \fg{}

\bigskip

\textbf{Q3.2 :} \og{} Quelles sont les limites techniques ou humaines que vous rencontrez avec ces m\'{e}thodes
actuelles (temps de r\'{e}action, faux positifs, charge de travail)~? \fg{}

\bigskip

\textbf{Q3.3 :} \og{} Face \`{a} la complexit\'{e} des attaques, que penseriez-vous de l'int\'{e}gration d'un mod\`{e}le
d'Intelligence Artificielle capable d'analyser le trafic en temps r\'{e}el pour d\'{e}tecter de
mani\`{e}re autonome ces comportements suspects~? \fg{}

\bigskip

\textbf{Q3.4 :} \og{} Pensez-vous qu'un mod\`{e}le IA serait pertinent non seulement pour d\'{e}tecter, mais aussi
pour intervenir directement sur l'infrastructure (ex~: bloquer une IP, modifier
dynamiquement l'acc\`{e}s) pour vous faire gagner du temps~? \fg{}

% ============================================================
%  4. PHASE DE CONCLUSION
% ============================================================
\section{Phase de Conclusion (Crit\`{e}res d'acceptabilit\'{e} de la solution)}

\textbf{Q4.1 :} \og{} Pour qu'une solution bas\'{e}e sur l'IA vous soit r\'{e}ellement utile, quelles devraient \^{e}tre
ses caract\'{e}ristiques principales (ex~: type de tableau de bord, alertes en temps r\'{e}el,
visibilit\'{e} du r\'{e}seau)~? \fg{}

\bigskip

\textbf{Q4.2 :} \og{} Quelles seraient vos exigences strictes ou vos craintes (ex~: confidentialit\'{e} des flux,
blocage accidentel d'un trafic l\'{e}gitime) avant d'autoriser une IA \`{a} interagir avec vos
routeurs/pare-feux~? \fg{}

\bigskip

\textbf{Q4.3 :} \og{} En r\'{e}sum\'{e}, si ce syst\`{e}me intelligent et autonome, garantissant un faible taux
d'erreur, \'{e}tait d\'{e}ploy\'{e} sur votre infrastructure, seriez-vous pr\^{e}t \`{a} l'adopter dans vos
op\'{e}rations quotidiennes~? \fg{}

\vspace{2em}
\noindent\rule{\linewidth}{0.4pt}
\begin{center}
  \small Guide d'Entretien --- Module Intelligence Artificielle --- Projet Acad\'{e}mique --- ENSA K\'{e}nitra
\end{center}

\end{document}